\section[On Encodings of the $\circuit_s(f)$ Tautology]{On Encodings
  of the $\bm{\mathrm{Circuit}_s(f)}$ Tautology}
\label{sec:encoding}

Let us introduce a possible constant width CNF encoding of
$\circuit_s(f)$ as proposed by Razborov \cite{Razborov98}.

The formula is defined over the same variables as introduced in
\cref{sec:circuit-encoding}, but in order to keep the fan-in bounded,
we further introduce the extension variables $\isvarless(v, a, i)$ and
$\isgateless(v, a, u)$ that indicate whether the wire $a$ of $v$ is
connected to a variable in $x_1, \ldots, x_i$, a gate $1, \ldots, u$
respectively.

Let us group the axioms in the same manner as we did in
\cref{sec:circuit-encoding}. First we have the structure axioms. The
first axioms encode that each wire is connected to a single kind
\begin{align}
  &\big(\isfromconst(v, a) \vee \isfromvar(v, a) \vee \isfromgate(v,
    a)\big) \wedge \nonumber\\
  &\neg\big(\isfromconst(v, a) \wedge \isfromvar(v, a)\big) \wedge\nonumber\\
  &\neg\big(\isfromvar(v, a) \wedge \isfromgate(v, a)\big) \wedge\nonumber\\
  &\neg\big(\isfromconst(v, a) \wedge \isfromgate(v, a)\big) \eqperiod
\end{align}
The next set of axioms similarly ensures that each gate computes
precisely one function
\begin{align}
  &\big(\isneg(v) \vee \isor(v) \vee \isand(v)\big) \land\nonumber\\
  &\neg\big(\isneg(v) \wedge \isor(v)\big) \wedge \neg\big(\isor(v) \wedge
    \isand(v)\big) \wedge \neg\big(\isneg(v) \wedge \isand(v)\big) \eqperiod
\end{align}
Last, we need to make sure that each wire is connected to a single
input or a gate.
\begin{align}
  \label{ax:input-cnf}
  &\isvarless(v, a, n) \wedge \bigwedge_{i \neq j}\neg \big(\isvar(v,a,i) \wedge
    \isvar(v,a,j)\big) \wedge\nonumber\\
    &\bigwedge_{i \in [n]}
      \Big(\isvarless(v, a, i) \equiv
    \big(\isvarless(v, a, i-1) \vee \isvar(v, a, i)\big)
    \Big) \eqcomma \nonumber\\
  &\qquad\text{where } \isvarless(v, a, 0) \defeq 0 \eqcomma
\end{align}
and similarly for $v \in [s] \setminus \{1\}$ we have that 
\begin{align}
  \label{ax:gate-cnf}
  &\isgateless(v, a, v-1) \wedge
    \bigwedge_{u < u' < v} \neg\big(\isgate(v,a,u) \wedge \isgate(v,a,u')\big)
    \wedge\nonumber\\
  &\bigwedge_{u \in [v-1]}
    \Big(
    \isgateless(v, a, u) \equiv
    \big(\isgateless(v, a, u-1) \vee \isgate(v, a, u)\big)
    \Big)\eqcomma \nonumber\\
  &\qquad\text{where }\isgateless(v, a, 0) \defeq 0 \eqperiod
\end{align}

Let us take a look at the evaluation axioms. Again, we have axioms
that ensure that the wires carry the values intended by the structure
variables. If a wire is connected to a constant, then the evaluation
variable associated with that wire should be equal to the constant
\begin{align}
  &\isfromconst(v, a)
    \rightarrow
    \big(\inwire_{
    \alpha}(v, a) \equiv \constval(v, a)\big) \eqcomma
\end{align}
and similarly if a wire is connected to an input or a gate
\begin{align}
  &\isfromvar(v, a) \wedge \isvar(v, a, i)
    \rightarrow
    \inwire_{\alpha}(v, a) \equiv \alpha_i \eqcomma\\
  &\isfromgate(v, a) \wedge \isgate(v, a, u)
    \rightarrow
    \inwire_{\alpha}(v, a) \equiv \outwire_{\alpha}(u) \eqperiod
\end{align}
Last we need to make sure that the gates propagate the value they are
supposed to compute.
\begin{align}
  &\isneg(v) \rightarrow
    \big(\outwire_{\alpha}(v) \equiv \neg \inwire_{\alpha}(v, 1)\big)\\
  &\isor(v) \rightarrow
    \big(\outwire_{\alpha}(v) \equiv \inwire_{\alpha}(v, 1) \vee \inwire_{\alpha}(v, 2)\big)\\
  &\isand(v) \rightarrow
    \big(\outwire_{\alpha}(v) \equiv \inwire_{\alpha}(v, 1) \wedge
    \inwire_{\alpha}(v, 2)\big)
    \eqperiod
\end{align}

The final axioms ensure that the correct function is computed
\begin{align}
  \outwire_{\alpha}(s) \equiv f(\alpha) \eqperiod
\end{align}

This formula can be rewritten in the usual manner into a $4$-CNF. Let
us denote this formula by $\circuit^{\mathrm{CNF}}_s(f)$.

\pa{Here the $4$ is because all axioms use at most 4 variables, right?  Can say this.}

\kr{It is a bit annoying: you have to consider it as a CNF and then
  each clause is over at most $4$ vars. The axioms themselves (for
  example \cref{ax:gate-cnf}) may depend on more variables. }

\begin{proposition}
  \label{prop:cnf-poly}
  If there is an SoS refutation of degree $d$ of the CNF formula
  $\circuit^{CNF}_s(f)$, then there is an SoS refutation of degree
  $O(d)$ of the system of polynomials $\circuit_s(f)$ as introduced in
  \cref{sec:circuit-encoding}.
\end{proposition}

The rest of this section is devoted to the proof of
\cref{prop:cnf-poly}.

Observe that for each axiom $p$ from the polynomial encoding
$\circuit_s(f)$, there is a CNF
$F_p \subseteq \circuit_s^{\mathrm{CNF}}(f)$ over the same variables
as $p$ (ignoring the added extension variables $\isvarless(v, a, i)$
and $\isgateless(v, a, u)$) such that $p(\alpha) = 0$ is satisfied by
a Boolean assignment $\alpha$ if and only if $F_p$ is satisfied by
$\alpha$ (where we extend the assignment to the extension variables in
the natural manner).

\pa{``extension variables'' not defined.  When introducing them, say
  that they are extension variables, and remind here what the
  extension variables are.  But I also don't understand what the
  purpose of this paragraph is, why is it needed?}

\kr{added the definition of extension variables and remind the reader
  by adding the names. The purpose of this paragraph is to introduce
  the notation $F_p$.}

Recall that SoS operates on polynomials and we thus need to translate
the CNF into a system of polynomials. We translate a clause
$\vee_{i \in [w]} z_i$ into the polynomial
$\prod_{i \in [w]}(1 - z_i) = 0$.

Observe that almost all axioms $p$ of $\circuit_s(f)$ depend only on a
constant number of variables. From such $p$, using the appropriate
Boolean axioms and negation axioms, we can in constant degree derive
$F_p$.

Let us define a polynomial substitution $\rho$ that gets rid of the
extension variables in the natural manner: the substitution $\rho$
first replace each occurrence of a ``bar'' extension variable
$\overline{\isvarless(v, a, i)}$ or $\overline{\isgateless(v, a, u)}$
by the polynomial $1-\isvarless(v, a, i)$ and the polynomial
$1-\isgateless(v, a, u)$ respectively. Then, $\rho$ replaces each
occurrence of the variable $\isvarless(v, a, i)$ by
$\sum_{j \le i} \isvar(v, a, j)$ and similarly $\isgateless(v, a, u)$
by $\sum_{w \le u} \isgate(v, a, w)$.

Suppose we have a degree $d$ refutation $\pi$ of
$\circuit_s^{\mathrm{CNF}}(f)$. Let us consider $\restrict{\pi}{\rho}$
and $\restrict{\circuit^{\mathrm{CNF}}(f)}{\rho}$. Note that
$\restrict{\pi}{\rho}$ is a degree $d$ SoS refutation of
$\restrict{\circuit^{\mathrm{CNF}}_s(f)}{\rho}$.

We claim that in constant degree the axioms of
$\restrict{\circuit^{\mathrm{CNF}}_s(f)}{\rho}$ can be derived from the
polynomial encoding $\circuit_s(f)$. As previously noted, this holds
for all axioms but the ones that are over a non-constant number of
variables. In other words it just remains to show that we can derive
the substituted \cref{ax:input-cnf,ax:gate-cnf} from
\cref{eq:ax-in-gt,eq:ax-in-gt2,eq:prod-zero-in,eq:prod-zero-gate}.

Let us consider \cref{ax:input-cnf}. With the extension variables
substituted and translated into a system of polynomials the axiom
consists of the following polynomial equations.
\begin{align}
  \label{eq:ax1}
  1 - \sum_{j \in [n]} \isvar(v,a,j) &= 0\\
  \label{eq:ax2}
  \isvar(v,a,i) \cdot \isvar(v,a,j) &= 0,
  \text{ for~}i \neq j\\
  \nonumber
  \left(\sum_{j \le i} \isvar(v,a,j)\right)
  \left(1 - \sum_{j < i} \isvar(v,a,j)\right)\cdot\quad&\\
  \label{eq:ax3}
  \overline{\isvar(v,a,i)} &= 0,
  \text{ for~}i \in [n]\\
  \label{eq:ax4}
  \left(1 - \sum_{j \le i} \isvar(v,a,j)\right)
  \left(\sum_{j < i} \isvar(v,a,j)\right) &= 0,
  \text{ for~}i \in [n]\\
  \label{eq:ax5}
  \left(1 - \sum_{j \le i} \isvar(v,a,j)\right)
  \isvar(v,a,i) &= 0,
  \text{ for~}i \in [n] \eqperiod
\end{align}
\Cref{eq:ax1} is equal to \cref{eq:ax-in-gt} and similarly
\cref{eq:ax2} is equal to \cref{eq:prod-zero-in}. In the following we
show that \cref{eq:ax3,eq:ax4,eq:ax5} can be derived from
\cref{eq:prod-zero-in}, the Boolean axioms and the negation axioms in
constant degree.

Consider \cref{eq:ax3}. Expand and rewrite modulo the Boolean axioms
and the negation axiom to obtain
% \begin{align}
%   \sum_{j \le i} \isvar(v,a,j) -
%   \isvar(v,a,i)\left(\sum_{j \le i} \isvar(v,a,j)\right) -\\
%   \left(\sum_{j \le i} \isvar(v,a,j)\right)
%   \left(\sum_{j < i} \isvar(v,a,j)\right) +\\
%   \isvar(v,a,i)
%   \left(\sum_{j \le i} \isvar(v,a,j)\right)
%   \left(\sum_{j < i} \isvar(v,a,j)\right)
%    = 0 \eqperiod
% \end{align}
% Rewrite modulo the Boolean axioms to get to
% \begin{align}
%   \sum_{j < i} \isvar(v,a,j) -
%   \isvar(v,a,i)\left(\sum_{j < i} \isvar(v,a,j)\right) -\\
%   \sum_{j < i} \isvar(v,a,j) -
%   2\sum_{j < j' < i}\isvar(v,a,j)\isvar(v,a,j') -
%   \isvar(v,a,i)\sum_{j < i}\isvar(v,a,j) +\\
%   \isvar(v,a,i)
%   \left(\sum_{j < i} \isvar(v,a,j)\right) +
%   \isvar(v,a,i)
%   \left(\sum_{j < i} \isvar(v,a,j)\right)^2
%   = 0 \eqperiod
% \end{align}
% We can cancel some terms and are left with the polynomial equation
\begin{align}
  &\isvar(v,a,i)
  \left(\sum_{j < i} \isvar(v,a,j)\right)^2-\nonumber\\
  &\qquad\qquad 2\sum_{j < j' < i}\isvar(v,a,j)\cdot\isvar(v,a,j')~-\nonumber\\
  &\qquad\qquad\qquad\qquad\isvar(v,a,i)\sum_{j < i}\isvar(v,a,j)
  = 0 \eqperiod
\end{align}
Observe that every term $t$ left in this polynomial is of the form
$t = t' \cdot \isvar(v,a,j) \cdot \isvar(v,a,j')$, for some
$j \neq j'\in [i]$ and a term $t'$ of degree at most $1$. But this
means that every term is equal to $0$ modulo \cref{eq:prod-zero-in}
and we thus see that \cref{eq:ax3} can be derived in constant degree
from $\circuit_s(f)$.

Let us consider \cref{eq:ax4}. Rewrite modulo the Boolean axiom to
obtain
\begin{align}
  &\isvar(v,a,i)\sum_{j < i}\isvar(v,a,j) +\nonumber\\
  &\qquad 2 \sum_{j < j' < i}\isvar(v,a,j)\cdot \isvar(v,a,j')
  = 0 \eqperiod
\end{align}
All terms are of the form of \cref{eq:prod-zero-in} and we can thus
derive \cref{eq:ax4} from $\circuit_s(f)$ in constant degree.

Last, we need to consider \cref{eq:ax5}. Note that modulo the Boolean
axiom we obtain the polynomial equation
\begin{align}
  - \isvar(v,a,i)\sum_{j < i}\isvar(v,a,j) = 0 \eqperiod
\end{align}
Also in this polynomial every term is of the form of
\cref{eq:prod-zero-in} and thus also \cref{eq:ax5} can be derived in
constant degree.

What remains is to show that \cref{ax:gate-cnf} can be derived from
$\circuit_s(f)$ in constant degree. This can be checked analogous to
\cref{ax:input-cnf} and we thus omit it here.

We conclude that all axioms of $\restrict{\circuit^{\mathrm{CNF}}_s(f)}{\rho}$ can
be derived from $\circuit_s(f)$ in constant degree and thus a degree
$d$ SoS refutation of $\circuit_s^{\mathrm{CNF}}(f)$ gives rise to a degree
$O(d)$ SoS refutation of $\circuit_s(f)$. Equivalently, a degree $d$
lower bound for $\circuit_s(f)$ implies a degree $\Omega(d)$ lower
bound for $\circuit_s^{\mathrm{CNF}}(f)$ as claimed.
